% Part: first-order-logic
% Chapter: syntax-and-semantics
% Section: unique-readability

\documentclass[../../../include/open-logic-section]{subfiles}

\begin{document}

\olfileid[pt]{fol}{syn}{unq}

\olsection{Leitura Única}

\begin{explain}
A maneira como definimos !!{formula}s garante que toda !!{formula} tem
uma \emph{leitura única}, isto é, há essencialmente apenas uma maneira de
construí-la de acordo com nossas regras de formação para !!{formula}s e
apenas uma maneira de ``interpretá-la''.  Se não fosse assim, teríamos
!!{formula}s ambíguas, isto é, !!{formula}s que têm mais de
uma leitura ou interpretação---e isso é claramente algo que queremos
evitar.  Mas, mais importante, sem essa propriedade, a maioria das
definições e provas que vamos dar não funcionaria.

Talvez a melhor maneira de deixar isso claro seja ver o que aconteceria se
tivéssemos dado regras ruins para formar !!{formula}s que não
garantissem a leitura única.  Por exemplo, poderíamos ter esquecido
os parênteses nas regras de formação para conectivos, por exemplo, poderíamos
ter permitido isto:
\begin{quote}
Se $!A$ e $!B$ são !!{formula}s, então $!A \lif !B$ também é.
\end{quote}
Partindo de uma fórmula atômica $!D$, isso nos permitiria formar $!D
\lif !D$. A partir disso, juntamente com $!D$, obteríamos $!D \lif !D
\lif !D$.  Mas há duas maneiras de fazer isso:
\begin{enumerate}
\item Tomamos $!D$ como $!A$ e $!D \lif !D$ como $!B$.
\item Tomamos $!A$ como $!D \lif !D$ e $!B$ é~$!D$.  
\end{enumerate}
Correspondentemente, há duas maneiras de
``ler'' a !!{formula}~$!D \lif !D \lif !D$. Ela é da forma $!B
\lif !C$ onde $!B$ é $!D$ e $!C$ é $!D \lif !D$, mas \emph{ela é
  também} da forma $!B \lif !C$ com $!B$ sendo $!D \lif !D$ e $!C$
sendo~$!D$.

Se isso acontece, nossas definições nem sempre funcionarão. Por exemplo,
quando definimos o !!{main operator} de uma fórmula, dizemos: em uma
fórmula da forma $!B \lif !C$, o !!{main operator} é a
ocorrência indicada de~$\lif$. Mas se podemos casar a fórmula $!D
\lif !D \lif !D$ com $!B \lif !C$ das duas maneiras diferentes mencionadas
acima, então em um caso obtemos a primeira ocorrência de $\lif$ como o
!!{main operator}, e no segundo caso a segunda ocorrência.  Mas
pretendemos que o !!{main operator} seja uma \emph{função} da
!!{formula}, isto é, toda !!{formula} deve ter exatamente uma ocorrência de
!!{main operator}.
\end{explain}

\begin{lem}
O número de parênteses esquerdos e direitos em !!a{formula}~$!A$ é
igual.
\end{lem}

\begin{proof}
Provamos isso por indução sobre a maneira como $!A$ é construída. Isso
requer duas coisas: (a) Temos primeiro que provar que todas as
fórmulas atômicas têm a propriedade em questão (a base da indução). (b) Então
temos que provar que, quando construímos novas fórmulas a partir de fórmulas dadas,
as novas fórmulas têm a propriedade, desde que as antigas a tenham.

Seja $l(!A)$ o número de parênteses esquerdos, e $r(!A)$ o número
de parênteses direitos em~$!A$, e $l(t)$ e $r(t)$ similarmente o
número de parênteses esquerdos e direitos em um termo~$t$.

\begin{prob}
    Prove que, para qualquer termo~$t$, $l(t) = r(t)$.
\end{prob}

\begin{enumerate}
\tagitem{prvFalse}{\indcase{!A}{\lfalse}{$\indfrm$ tem $0$ parênteses esquerdos e $0$
    parênteses direitos.}}{}

\tagitem{prvTrue}{\indcase{!A}{\ltrue}{$\indfrm$ tem $0$ parênteses esquerdos e $0$
    parênteses direitos.}}{}

\item \indcase{!A}{\Atom{R}{t_1,\dots,t_n}}{$l(\indfrm) = 1 + l(t_1) +
  \dots + l(t_n) = 1 + r(t_1) + \dots + r(t_n) = r(\indfrm)$. Aqui
  usamos o fato, deixado como exercício, de que $l(t) = r(t)$ para
  qualquer termo~$t$.}

\item \indcase{!A}{\eq[t_1][t_2]}{$l(\indfrm) = l(t_1) + l(t_2) =
  r(t_1) + r(t_2) = r(\indfrm)$.}

\tagitem{prvNot}{\indcase{!A}{\lnot !B}{Por hipótese de indução,
    $l(!B) = r(!B)$. Assim $l(\indfrm) = l(!B) = r(!B) =
    r(\indfrm)$.}}{}

\item \indcase{!A}{(!B \ast !C)}{Por hipótese de indução, $l(!B) =
  r(!B)$ e $l(!C) = r(!C)$.  Assim $l(\indfrm) = 1 + l(!B) + l(!C) =
  1 + r(!B) + r(!C) = r(\indfrm)$.}

\tagitem{prvAll}{\indcase{!A}{\lforall[x][!B]}{Por hipótese de
    indução, $l(!B) = r(!B)$. Assim, $l(\indfrm) = l(!B) = r(!B) =
    r(\indfrm)$.}}{}

% Print case for \lexists if prvEx and not prvAll
\tagitem{prvAll,notprvEx}{}{\indcase{!A}{\lexists[x][!B]}{Por hipótese de
    indução, $l(!B) = r(!B)$. Assim, $l(\indfrm) = l(!B) = r(!B) =
    r(\indfrm)$.}}

% Just say similarly if prvEx and prvAll
\tagitem{notprvAll,notprvEx}{}{\indcase{!A}{\lexists[x][!B]}{Analogamente.}}
\end{enumerate}
\end{proof}

\begin{defn}[Prefixo próprio]
Uma cadeia de símbolos $!B$ é um \emph{prefixo próprio} de uma cadeia de símbolos~$!A$ se
a concatenação de $!B$ e uma cadeia de símbolos não vazia resulta em~$!A$.
\end{defn}

\begin{lem}\ollabel{lem:no-prefix}
Se $!A$ é !!a{formula}, e $!B$ é um prefixo próprio de $!A$, então
$!B$ não é !!a{formula}.
\end{lem}

\begin{proof}
Exercício.
\end{proof}

\begin{prob}
Prove \olref[fol][syn][unq]{lem:no-prefix}.
\end{prob}

\begin{prop}
\ollabel{prop:unique-atomic}
Se $!A$ é uma !!{formula} atômica, então ela satisfaz uma, e apenas uma,
das seguintes condições.
\begin{enumerate}
\tagitem{prvFalse}{$!A \ident \lfalse$.}{}
\tagitem{prvTrue}{$!A \ident \ltrue$.}{}
\item $!A \ident \Atom{R}{t_1,\dots,t_n}$ onde $R$ é um !!{predicate}
  de $n$ lugares, $t_1$, \dots, $t_n$ são termos, e cada um de $R$,
  $t_1$, \dots, $t_n$ é determinado de forma única.
\item $!A \ident \eq[t_1][t_2]$ onde $t_1$ e $t_2$ são termos
  determinados de forma única.
\end{enumerate}
\end{prop}

\begin{proof}
Exercício.
\end{proof}

\begin{prob}
Prove \olref[fol][syn][unq]{prop:unique-atomic} (Dica: Formule e
prove uma versão de \olref[fol][syn][unq]{lem:no-prefix} para termos.)
\end{prob}

\begin{prop}[Leitura Única]
Toda !!{formula} satisfaz uma, e apenas uma, das seguintes condições.
\begin{enumerate}
\item $!A$ é atômica.

\tagitem{prvNot}{$!A$ é da forma $\lnot !B$.}{}

\tagitem{prvAnd}{$!A$ é da forma $(!B \land !C)$.}{}

\tagitem{prvOr}{$!A$ é da forma $(!B \lor !C)$.}{}

\tagitem{prvIf}{$!A$ é da forma $(!B \lif !C)$.}{}

\tagitem{prvIff}{$!A$ é da forma $(!B \liff !C)$.}{}

\tagitem{prvAll}{$!A$ é da forma $\lforall[x][!B]$.}{}

\tagitem{prvEx}{$!A$ é da forma $\lexists[x][!B]$.}{}
\end{enumerate}
Além disso, em cada caso $!B$, ou $!B$ e $!C$, são determinados de forma
única.  Isso significa que, por exemplo, não há pares diferentes $!B$,
$!C$ e $!B'$, $!C'$ tais que $!A$ seja ao mesmo tempo da forma 
\iftag{prvIf}{$(!B \lif !C)$ e $(!B' \lif !C')$.}{
    \iftag{prvOr}{$(!B \lor !C)$ e $(!B' \lor !C')$.}{
      \iftag{prvAnd}{$(!B \land !C)$ e $(!B' \land !C')$.}{}}}
\end{prop}

\begin{proof}
As regras de formação exigem que, se !!a{formula} não é atômica, ela
deve começar com um parêntese de abertura~(, \iftag{prvNot}{$\lnot$,}{}
ou um quantificador. Por outro lado, toda !!{formula} que começa com
um dos seguintes símbolos deve ser atômica: !!a{predicate}, !!a{function}, !!a{constant}\iftag{prvFalse}{, $\lfalse$}{}\iftag{prvTrue}{, $\ltrue$}{}.

Então, na verdade só temos que mostrar que se $!A$ é da forma $(!B \ast
!C)$ e também da forma $(!B' \mathbin{\ast'} !C')$, então $!B \ident
!B'$, $!C \ident !C'$, e $\ast = {\ast'}$.

Então suponha que tanto $!A \ident (!B \ast !C)$ quanto $!A \ident (!B'
\mathbin{\ast'} !C')$.  Então ou $!B \ident !B'$ ou não.  Se for,
claramente $\ast = {\ast'}$ e $!C \ident !C'$, pois elas então são
subcadeias de $!A$ que começam no mesmo lugar e têm o mesmo
comprimento.  O outro caso é $!B \not\ident !B'$.  Como $!B$ e
$!B'$ são ambas subcadeias de $!A$ que começam no mesmo lugar, uma
deve ser um prefixo próprio da outra.  Mas isso é impossível por
\olref{lem:no-prefix}.
\end{proof}


\end{document}
